\documentclass[11pt,a4paper]{article}
\usepackage[T1]{fontenc}
\usepackage[utf8]{inputenc}
\usepackage{lmodern}
\usepackage[a4paper,margin=1in]{geometry}
\usepackage{setspace}
\onehalfspacing
\usepackage{amsmath,amssymb,amsthm,mathtools}
\usepackage{graphicx}
\graphicspath{{images/}}
\usepackage{subcaption}
\usepackage{booktabs}
\usepackage{array}
\usepackage{float}
\usepackage{xcolor}
\usepackage{siunitx}
\usepackage{enumitem}
\usepackage{csquotes}
\usepackage{hyperref}
\usepackage[nameinlink,noabbrev]{cleveref}

\hypersetup{
    colorlinks=true,
    linkcolor=blue!60!black,
    citecolor=green!50!black,
    urlcolor=magenta!60!black,
    pdftitle={Motion Planning PNO Documentation},
    pdfauthor={Your Name}
}

\newtheorem{theorem}{Theorem}
\newtheorem{lemma}{Lemma}
\newtheorem{proposition}{Proposition}
\theoremstyle{definition}
\newtheorem{definition}{Definition}
\newtheorem{remark}{Remark}

\title{Motion Planning with PNO\\Documentation}
\author{Samyak Soni}

\begin{document}

\maketitle
\tableofcontents
\newpage

\section{Problem Statement}\label{sec:problem}

\subsection{Objective}
Learn to compute a value function $V(x, g)$ from any 2D obstacle environment and goal position $(g)$.
The value function provides the minimum cost-to-goal from every spatial location $x$.

\subsection{Inputs}
\begin{itemize}
  \item \textbf{Binary Occupancy Map}: $m(x) \in \{0,1\}$ where $1$ indicates free space, $0$ indicates obstacles.
  \item \textbf{Goal Position}: $g = (g_x, g_y)$ as discrete pixel coordinates.
\end{itemize}

\subsection{Outputs}
\begin{itemize}
  \item \textbf{Value Function Field}: $V(x, g)$ where each pixel holds the shortest distance/cost to reach goal $g$ from that location.
  \item \textbf{Implicit Path}: The actual path is recovered via gradient descent on $V$: $\mathbf{p}(t+1) = \mathbf{p}(t) - \eta \nabla V(\mathbf{p}(t), g)$.
\end{itemize}

\subsection{Planning Extension}
The learned value function is designed to serve as an $\epsilon$-consistent heuristic for A* search.
This guarantees:
\begin{itemize}
  \item 100\% success rates (no goal-unreachability errors).
  \item Minimal node expansions via admissible heuristic guidance.
  \item Optimality with $\epsilon$-bounded cost approximation.
\end{itemize}

\section{Mathematical Formulation}\label{sec:math}

\subsection{The Eikonal Partial Differential Equation}
The value function satisfies the Eikonal PDE:
\[
  \|\nabla V(x, g)\| = c(x), \quad V(g, g) = 0
\]
where:
\begin{itemize}
  \item $\nabla V(x, g)$ is the spatial gradient of the value function.
  \item $c(x) = m(x)$ is the local traversal cost (1 in free space, $\infty$ in obstacles).
  \item The boundary condition $V(g, g) = 0$ anchors the value at the goal.
\end{itemize}

\subsection{Interpretation}
\begin{itemize}
  \item The magnitude of the gradient directly encodes the cost per unit distance.
  \item In free space, $|\nabla V| = 1$ implies Euclidean distance. 
  \item On or near obstacles, the gradient is undefined (non-smooth boundary).
  \item Shortest paths are found via $\nabla V$: steepest descent curves from any $(x)$ to $g$.
\end{itemize}

\section{System Pipeline}\label{sec:pipeline}

\noindent
The overall approach:
\begin{enumerate}
  \item \textbf{Input}: Binary occupancy map $m$ and goal $g$.
  \item \textbf{SDF Computation}: Convert binary map to Signed Distance Field (via FNO-based SDF model).
  \item \textbf{Planning Inference}: Use PNO to map (map, SDF, goal) to value function $V$.
  \item \textbf{Path Recovery}: Gradient descent on $V$ traces shortest path; A* integration uses $V$ as heuristic.
\end{enumerate}

\section{Model Architectures}\label{sec:architecture}

\subsection{Phase 1: SDF-FNO (Geometry to Signed Distance Field)}

\subsubsection{Motivation}
Binary occupancy maps are non-differentiable at obstacle boundaries. To enable smooth gradient-based operations,
we first convert the binary map to a continuous Signed Distance Field (SDF) using a Fourier Neural Operator (FNO).

\subsubsection{Architecture Specification}
The FNO2dSDF model processes single-channel input (binary map) to single-channel output (SDF):

\noindent
\textbf{Layer Structure:}
\begin{itemize}
  \item \textbf{Input Lift}: Linear projection from 1 channel $\to$ 32 channels (via \texttt{fc0}).
  \item \textbf{Spectral Blocks} ($\times 4$): Each block contains:
    \begin{itemize}
      \item \texttt{SpectralConv2d}: Operates in Fourier space; learns complex weights for the low $\text{modes}\_1 \times \text{modes}\_2$ frequencies.
      \item \texttt{Conv2d(1x1)}: Spatial mixing layer; refines feature interactions locally.
      \item Activation: GELU.
    \end{itemize}
  \item \textbf{Output Head}: Sequential MLP (32 $\to$ 128 $\to$ 1) with GELU; outputs single-channel SDF.
\end{itemize}

\subsection{Phase 2: Planning Neural Operator (PNO)}

\subsubsection{Core Innovation}
The PNO learns the Eikonal solution directly by jointly conditioning on three inputs:
the raw binary map, the continuous SDF, and the goal location. 
The PNO employs domain-awareness via the Smoothed Indicator Function (SIFN) mask to robustly handle obstacle boundaries.

\subsubsection{Architecture Specification}

\subsubsection{Input Construction}

The PNO concatenates three channels:
\[
  \mathbf{x}_{\text{in}} = [m(x), \text{SDF}(x), \mathbf{g}(x)] \in \mathbb{R}^{B \times 3 \times H \times W}
\]
where:
\begin{itemize}
  \item $m(x)$: Binary occupancy map.
  \item $\text{SDF}(x)$: Pre-computed signed distance field from Phase 1.
  \item $\mathbf{g}(x)$: Goal channel; a sparse tensor with $-1$ at goal coordinates, $0$ elsewhere.
\end{itemize}

\subsubsection{Smoothed Indicator Function (SIFN) Mask}

Before processing, a domain-aware mask is constructed using a FNO model:
\[
  \chi(x) = \sigma(\text{FNO}_{\text{mask}}(m, \text{SDF}))
\]
where $\sigma$ is a sigmoid activation ensuring $\chi \in (0,1)$.

\noindent
\textbf{Purpose:}
\begin{itemize}
  \item Smooth transition from obstacles ($\chi \approx 0$) to free space ($\chi \approx 1$).
  \item Allows spectral layers to learn localized boundary conditions.
  \item Prevents spurious oscillations near walls via soft masking.
\end{itemize}

\subsubsection{DAFNO Blocks (Domain-Agnostic Fourier Neural Operator Blocks)}

The core processing pipeline uses $4$ stacked DAFNO blocks. Each block:

\begin{itemize}
  \item \textbf{Spectral Layer (DAFNOSpectralConv2d)}:
    \begin{itemize}
      \item Computes 2D FFT of the masked feature: $\tilde{v} = \mathcal{F}(v \odot \chi)$.
      \item Applies learnable complex weights to low-frequency modes: $\tilde{v}_{\text{out}} = W \cdot \tilde{v}_{\text{in}}$.
      \item Inverses with 2D iFFT: $v_{\text{out}} = \mathcal{F}^{-1}(\tilde{v}_{\text{out}})$.
    \end{itemize}
  \item \textbf{Local Convolution}: 1$\times$1 Conv2d; enables point-wise non-linear feature mixing.
  \item \textbf{Residual Connection \& Activation}: 
    \[
      v_{\text{out}} = \text{GELU}(v_{\text{spec}} + v_{\text{local}})
    \]
\end{itemize}

\subsubsection{DeepNorm Projection Head}

The final layer enforces the Triangle Inequality constraint (value function consistency):
\[
  V(x, g) \leq V(x, x') + V(x', g) \quad \forall x, x', g
\]

\noindent
\textbf{Implementation:}
\begin{itemize}
  \item Computes feature differences relative to the goal: $\Delta v = |v - v_{\text{goal}}|$.
  \item Projects through a 2-layer MLP with \textbf{positive weights only} (enforced via \texttt{softplus}):
    \[
      V(x, g) = \sigma_2(W_2^+ \cdot \text{GELU}(W_1^+ \cdot \Delta v + b_1) + b_2)
    \]
    where $W_i^+ = \text{softplus}(W_i)$.
  \item Positivity of weights ensures monotonic increase in cost with spatial distance, preserving optimality.
\end{itemize}


\end{document}
